\documentclass{article}

\usepackage{fancyhdr}
\usepackage{amsmath}
\usepackage{amsthm}
\usepackage{amsfonts}
\usepackage{tikz}
\usepackage[plain]{algorithm}
\usepackage{algpseudocode}
\usepackage{amssymb}
\usepackage{enumitem}

\usetikzlibrary{automata,positioning}

\topmargin=-0.45in
\evensidemargin=0in
\oddsidemargin=0in
\textwidth=6.5in
\textheight=9.0in
\headsep=0.25in

\linespread{1.1}

\pagestyle{fancy}
\lhead{\hmwkAuthorName}
\chead{}
\rhead{\hmwkHeaderTitle}
\cfoot{\thepage}

\renewcommand\headrulewidth{0.4pt}
\renewcommand\footrulewidth{0pt}

\setlength\parindent{0pt}

\newcommand{\hmwkTitle}{Homework\ 8}
\newcommand{\hmwkDueDate}{Monday, March 30, 2026,\ 11{:}59 pm (Thai time)}
\newcommand{\hmwkClass}{Data Structure\ and Abstractions}
\newcommand{\hmwkClassTime}{Section 2}
\newcommand{\hmwkClassInstructor}{Kanat Tangwongsan}
\newcommand{\hmwkAuthorName}{\textbf{Jiraroj Wiruchpongsanon (6781617)}}

\newcommand{\hmwkHeaderTitle}{Data Structure and Abstractions: Homework 8}

\title{
    \vspace{2in}
    \textmd{\textbf{\hmwkClass:\ \hmwkTitle}}\\
    \normalsize\vspace{0.1in}\small{Due\ on\ \hmwkDueDate}\\
    \vspace{0.1in}\large{\textit{\hmwkClassInstructor\ \hmwkClassTime}}
    \vspace{3in}
}

\author{\hmwkAuthorName}
\date{}

\renewcommand{\part}[1]{\textbf{\large Part \Alph{partCounter}}\stepcounter{partCounter}\\}

\newcounter{partCounter}
\newcommand{\solution}{\textbf{\large Solution}}

\newtheorem*{theorem}{Theorem}

\theoremstyle{remark}
\newtheorem*{remark}{Remark}

\begin{document}

\maketitle
\begin{center}
    \textit{Data Struct. and \{Abstractions, OOP\} (Term II/2025--26)}\\
    \textit{built on 2026/03/19 at 13:42:12}\\
    \textit{due: mon mar 30 @ 11:59pm}
\end{center}
\newpage

This assignment contains both programming and written questions. For the written portion, you
will typeset your answer, produce a PDF file called \texttt{hw8.pdf}, and upload it to Canvas. No other format
will be accepted.

\subsection*{Collaboration}

We interpret collaboration very liberally. You may work with other students. However, each student
must write up and hand in his or her assignment separately. Let us repeat: You need to write your
own code. You must not look at or copy someone else’s code. You need to write up answers to written
problems individually. The fact that you can recreate the solution from memory will be taken as proof
that you actually understood it, and you may actually be interviewed about your answers.

\newpage

\section*{Task 1: Breadth-First Search Running Time (2 points)}

Our breadth-first (BFS) implementation from class uses HashSets and HashMaps to keep track of
progress. In this task, you’ll reanalyze the running time of BFS if we use TreeSets and TreeMaps instead.
Show your work carefully. To keep things simple, we’ll work with the BFS code shown above. It only
reports vertices reachable from a source s. It has been rewritten from the version presented in class to
make the operations more explicit.

\begin{verbatim}
Set<Integer> nbrsExcluding(
    UndirectedGraph<Integer> G,
    Set<Integer> vtxes,
    Set<Integer> excl
) {
    Set<Integer> union = new TreeSet<>(); // not HashMap
    for (Integer src : vtxes) {
        for (Integer dst : G.adj(src))
            if (!excl.contains(dst)) { union.add(dst); }
    }
    return union;
}

Set<Integer> bfs(UndirectedGraph<Integer> G, int s) {
    Set<Integer> frontier = new TreeSet<>(Arrays.asList(s));
    Set<Integer> visited = new TreeSet<>(Arrays.asList(s));
    while (!frontier.isEmpty()) {
        frontier = nbrsExcluding(G, frontier, visited);
        visited.addAll(frontier); // the i-th position is what's reached at i hops
    }
    return visited;
}
\end{verbatim}

You should know that the TreeSet implementation uses a balanced binary search tree (BST), so
add, contains, and remove, unlike in a HashSet, take O(logS) per operation, where S is the size of the
collection. Though addAll may perform other optimizations, for the purpose of this task, assume that
addAll(X) simply repeatedly calls add on each element of X .

Homework 8 Data Struct.
You should also assume that UndirectedGraph G is implemented as an adjacency table using
HashMaps, so G.adj takes O(1).

\medskip
\solution

\vspace{1em}
\hrulefill

\section*{Task 2: Random Permutations (2 points)}

Remember that [n] is the set \{1,2,...,n\}. Let p : [n] $\to$ [n] be a random permutation on [n] chosen
uniformly among the n! permutations. Consider the following code:

\begin{verbatim}
int minSoFar = Integer.MAX_VALUE;
int numUpdate = 0;
for (int i=1;i<=n;i++) {
    if (p(i) < minSoFar) {
        minSoFar = p(i);
        numUpdate++;
    }
}
\end{verbatim}

We’ll see that numUpdate is a random variable that depends on the permutation randomly chosen.
Prove that at the end of the for-loop,

\[\ln(n +1) \le \mathbf{E}[\text{numUpdate}] \le 1+\ln n\]

\vspace{1em}
\hrulefill

% hackerrank id

\end{document}
